\documentclass[11pt]{article}

\usepackage{amsmath, amssymb}
\usepackage{graphicx}
\usepackage{hyperref}
\usepackage{geometry}
\geometry{margin=1in}

\title{Sovereign Inference-as-a-Service:\\
Decentralized Discovery, Authentication, and Execution of AI Inference}
\author{Jean-François Bouchard}
\date{April 2026}

\begin{document}

\maketitle

\begin{abstract}
Contemporary AI inference services rely on centralized infrastructure, including API gateways, revocable credentials, and network endpoints subject to filtering or seizure.

We present a sovereign inference-as-a-service architecture that removes reliance on centralized control points. Service discovery is performed via SRV records stored in the Emercoin Name-Value Storage (NVS) blockchain. Trust is established through on-chain identities (privatenessID), replacing certificate authority dependencies with cryptographic key bindings anchored in a public ledger.

Transport is provided by mesh and overlay networks, reducing dependence on conventional IP routing visibility. Payment is performed using native blockchain transactions without intermediary processors.

While network-level interference remains possible, the system ensures that service discovery and identity records remain globally replicated and tamper-resistant.

We further introduce the use of a user-controlled entropy source for inference sampling, reducing the risk of operator-induced bias in probabilistic generation.

The architecture is validated on commodity hardware, establishing a reproducible lower bound for decentralized inference deployment.
\end{abstract}

\section{Introduction}

AI inference delivery is predominantly centralized, relying on cloud infrastructure and revocable access mechanisms. This introduces structural dependencies on infrastructure providers, payment systems, and regulatory environments.

This paper proposes a decentralized alternative in which:
\begin{itemize}
    \item Discovery is blockchain-based
    \item Identity is cryptographically anchored
    \item Transport is multi-network
    \item Payment is native and decentralized
\end{itemize}

\section{System Model}

We define the system as:

\[
\mathcal{S} = (N, I, D, T, E)
\]

Where:
\begin{itemize}
    \item $N$: set of nodes
    \item $I$: identity mapping (on-chain)
    \item $D$: discovery mechanism
    \item $T$: transport layer
    \item $E$: entropy source
\end{itemize}

\section{Consensus and Storage}

All persistent records are stored in Emercoin NVS. Records are:
\begin{itemize}
    \item globally replicated
    \item cryptographically committed
    \item modifiable only with key ownership
\end{itemize}

Auxiliary proof-of-work (AuxPoW) increases resistance to deep chain reorganization by coupling security to external mining participation.

\section{Security Properties}

\subsection{Record Immutability}

Unauthorized modification requires:

\[
P(\text{attack}) \approx P(\text{key compromise}) + P(\text{consensus failure})
\]

\subsection{Discovery Integrity}

A client accepts an endpoint only if:

\[
\text{Endpoint Key} \in I(\text{identity})
\]

\subsection{No Single Point of Control}

\[
\forall c \in \mathcal{S}, \quad \mathcal{S} \setminus c \neq \emptyset
\]

\subsection{Attack Cost Lower Bound}

\[
C_{attack} \ge \min(C_{reorg}, C_{key}) + C_{network}
\]

\section{Entropy as a Security Primitive}

\subsection{Sampling Model}

\[
P(t_i) = \frac{\exp(l_i / T)}{\sum_j \exp(l_j / T)}
\]

\subsection{Adversarial Bias}

\[
P'(t_i) = P(t_i) + \epsilon_i
\]

Where:

\[
\sum_i \epsilon_i = 0
\]

\subsection{Detection Difficulty}

\[
D_{KL}(P' || P) \approx \mathcal{O}(\epsilon^2)
\]

Small biases are difficult to detect without large datasets.

\subsection{Mitigation}

External entropy source:

\[
R = H(seed, counter)
\]

Sampling becomes:

\[
t_i = f(P, R)
\]

Reducing operator control over stochastic output.

\section{Experimental Evaluation}

\subsection{Setup}

\begin{itemize}
    \item Model: 7B / 8B quantized
    \item Runtime: llama.cpp
    \item Hardware: Raspberry Pi 4, RK35xx
    \item Transports: localhost, LAN, mesh, inet
\end{itemize}

\subsection{Throughput}

\begin{tabular}{|c|c|}
\hline
Hardware & Tokens/sec \\
\hline
Pi 4 & 4--8 \\
RK35xx & 8--18 \\
\hline
\end{tabular}

\subsection{Latency}

\begin{tabular}{|c|c|}
\hline
Transport & Latency \\
\hline
Localhost & 50--150 ms \\
LAN & 100--300 ms \\
Mesh & 500--1500 ms \\
Inet & 200--800 ms \\
\hline
\end{tabular}

\subsection{Discovery Time}

\begin{tabular}{|c|c|}
\hline
Operation & Time \\
\hline
NVS lookup & 1--3 s \\
SRV resolution & 50--200 ms \\
\hline
\end{tabular}

\section{Economic Considerations}

Operational costs are low due to:
\begin{itemize}
    \item commodity hardware
    \item minimal storage fees
\end{itemize}

Staking rewards may offset costs under favorable conditions, but this is dependent on market dynamics and not assumed as a system invariant.

\section{Limitations}

\begin{itemize}
    \item Network-level filtering remains possible
    \item Latency is higher on mesh transports
    \item Reputation systems are not yet implemented
    \item Real-time load balancing is limited
\end{itemize}

\section{Conclusion}

We present a decentralized inference architecture that:
\begin{itemize}
    \item removes reliance on centralized trust anchors
    \item distributes discovery and identity
    \item introduces entropy as a security control
\end{itemize}

The system demonstrates that AI inference can be delivered through decentralized infrastructure with verifiable trust and minimal cost.

\end{document}